\textbf{Phần 1. Trắc nghiệm lựa chọn}
\setcounterpageref{ex}{0}

\begin{ex} %[2H5N26]
Đường gấp khúc $ABCD$ gồm mấy đoạn thẳng?
\choice
{$2$ đoạn thẳng}
{\True $3$ đoạn thẳng ($AB, BC, CD$)}
{$4$ đoạn thẳng}
{$5$ đoạn thẳng}
\loigiai{
Gồm 3 đoạn thẳng $AB, BC, CD$.\\
\textbf{Đáp án của câu hỏi là B}
}
\end{ex}

\begin{ex} %[2H5N26]
Hình tứ giác là hình có
\choice
{$3$ cạnh}
{$5$ cạnh}
{$6$ cạnh}
{\True $4$ cạnh}
\loigiai{
Tứ giác có 4 cạnh.\\
\textbf{Đáp án của câu hỏi là D}
}
\end{ex}

\begin{ex} %[2H5H26]
Độ dài đường gấp khúc $ABC$ có $AB = 4\text{ cm}, BC = 5\text{ cm}$ là
\choice
{$8\text{ cm}$}
{$10\text{ cm}$}
{\True $9\text{ cm}$}
{$1\text{ cm}$}
\loigiai{
$4 + 5 = 9\text{ cm}$.\\
\textbf{Đáp án của câu hỏi là C}
}
\end{ex}

\begin{ex} %[2H5H26]
Đường gấp khúc $MNPQ$ gồm các đoạn thẳng có độ dài lần lượt là $3\text{ cm}, 4\text{ cm}, 5\text{ cm}$. Độ dài đường gấp khúc đó là
\choice
{\True $12\text{ cm}$}
{$11\text{ cm}$}
{$13\text{ cm}$}
{$14\text{ cm}$}
\loigiai{
$3 + 4 + 5 = 12\text{ cm}$.\\
\textbf{Đáp án của câu hỏi là A}
}
\end{ex}

\begin{ex} %[2H5H26]
Hình chữ nhật và hình vuông đều là
\choice
{Hình tam giác}
{Đường gấp khúc}
{Hình tròn}
{\True Hình tứ giác}
\loigiai{
Hình chữ nhật, hình vuông đều là tứ giác có 4 cạnh.\\
\textbf{Đáp án của câu hỏi là D}
}
\end{ex}

\begin{ex} %[2H5H26]
Một hình tứ giác có độ dài các cạnh là $3\text{ cm}, 4\text{ cm}, 5\text{ cm}, 6\text{ cm}$. Tổng độ dài các cạnh của tứ giác là
\choice
{$17\text{ cm}$}
{\True $18\text{ cm}$}
{$19\text{ cm}$}
{$20\text{ cm}$}
\loigiai{
$3 + 4 + 5 + 6 = 18\text{ cm}$.\\
\textbf{Đáp án của câu hỏi là B}
}
\end{ex}

\begin{ex} %[2H5V26]
Đường gấp khúc $ABCD$ dài $24\text{ cm}$. Biết tổng độ dài đoạn $AB$ và $BC$ là $15\text{ cm}$. Độ dài đoạn $CD$ là
\choice
{\True $9\text{ cm}$}
{$8\text{ cm}$}
{$10\text{ cm}$}
{$11\text{ cm}$}
\loigiai{
$CD = 24 - 15 = 9\text{ cm}$.\\
\textbf{Đáp án của câu hỏi là A}
}
\end{ex}

\begin{ex} %[2H5V26]
Đường gấp khúc gồm 4 đoạn thẳng có độ dài bằng nhau, mỗi đoạn dài $5\text{ cm}$. Độ dài đường gấp khúc đó là
\choice
{$15\text{ cm}$}
{$25\text{ cm}$}
{$10\text{ cm}$}
{\True $20\text{ cm}$}
\loigiai{
$5 \times 4 = 20\text{ cm}$.\\
\textbf{Đáp án của câu hỏi là D}
}
\end{ex}

\begin{ex} %[2H5V26]
Một hình tứ giác có 4 cạnh bằng nhau, mỗi cạnh dài $6\text{ cm}$. Tổng độ dài các cạnh của hình tứ giác là
\choice
{$20\text{ cm}$}
{$28\text{ cm}$}
{\True $24\text{ cm}$}
{$18\text{ cm}$}
\loigiai{
$6 \times 4 = 24\text{ cm}$.\\
\textbf{Đáp án của câu hỏi là C}
}
\end{ex}

\begin{ex} %[2H5V26]
Một đường gấp khúc gồm 2 đoạn thẳng dài $18\text{ cm}$. Biết đoạn thứ nhất dài $8\text{ cm}$. Đoạn thứ hai dài hơn đoạn thứ nhất là
\choice
{\True $2\text{ cm}$}
{$1\text{ cm}$}
{$3\text{ cm}$}
{$4\text{ cm}$}
\loigiai{
Đoạn thứ hai: $18 - 8 = 10\text{ cm}$.\\
Đoạn 2 hơn đoạn 1: $10 - 8 = 2\text{ cm}$.\\
\textbf{Đáp án của câu hỏi là A}
}
\end{ex}

\begin{ex} %[2H5V26]
Độ dài đường gấp khúc $ABC$ là $16\text{ cm}$, biết đoạn thẳng $AB = BC$. Độ dài mỗi đoạn thẳng là
\choice
{$7\text{ cm}$}
{$9\text{ cm}$}
{\True $8\text{ cm}$}
{$6\text{ cm}$}
\loigiai{
$16 : 2 = 8\text{ cm}$.\\
\textbf{Đáp án của câu hỏi là C}
}
\end{ex}

\begin{ex} %[2H5C26]
Một đường gấp khúc khép kín gồm 4 đoạn thẳng tạo thành hình gì?
\choice
{Hình tam giác}
{\True Hình tứ giác}
{Đoạn thẳng}
{Hình tròn}
\loigiai{
Đường gấp khúc khép kín 4 đoạn tạo thành hình tứ giác.\\
\textbf{Đáp án của câu hỏi là B}
}
\end{ex}

\textbf{Phần 2. Trắc nghiệm đúng sai}
\setcounterpageref{ex}{0}

\begin{ex} %[2H5H26]
Xét tính đúng sai về đường gấp khúc và hình tứ giác:
\choiceTF
{\True Độ dài đường gấp khúc bằng tổng độ dài các đoạn thẳng tạo nên nó.}
{\True Hình tứ giác luôn có 4 cạnh.}
{Hình tam giác có 4 cạnh.}
{\True Hình tam giác có 3 cạnh, hình tứ giác có 4 cạnh.}
\loigiai{
\begin{itemize}
    \item Độ dài bằng tổng các đoạn: Đúng.
    \item Tứ giác có 4 cạnh: Đúng.
    \item Tam giác có 3 cạnh (không phải 4 cạnh). Ý 3 Sai, Ý 4 Đúng.
\end{itemize}
\begin{itemize}
    \item a) \textbf{Đúng}.
    \item b) \textbf{Đúng}.
    \item c) \textbf{Sai}.
    \item d) \textbf{Đúng}.
\end{itemize}
}
\end{ex}

\begin{ex} %[2H5V26]
Đường gấp khúc $MNP$ có $MN = 12\text{ cm}, NP = 15\text{ cm}$. Xét tính đúng sai:
\choiceTF
{\True Độ dài đường gấp khúc $MNP$ là $27\text{ cm}$ ($12 + 15 = 27$).}
{\True Đoạn $NP$ dài hơn đoạn $MN$ là $3\text{ cm}$ ($15 - 12 = 3$).}
{Đường gấp khúc $MNP$ gồm 3 đoạn thẳng.}
{\True Đường gấp khúc $MNP$ chỉ gồm 2 đoạn thẳng là $MN$ và $NP$.}
\loigiai{
\begin{itemize}
    \item Độ dài: $12 + 15 = 27\text{ cm}$: Đúng.
    \item $NP$ hơn $MN$ là $3\text{ cm}$: Đúng.
    \item Gồm 2 đoạn thẳng $MN$ và $NP$. Ý 3 Sai, Ý 4 Đúng.
\end{itemize}
\begin{itemize}
    \item a) \textbf{Đúng}.
    \item b) \textbf{Đúng}.
    \item c) \textbf{Sai}.
    \item d) \textbf{Đúng}.
\end{itemize}
}
\end{ex}

\textbf{Phần 3. Trả lời ngắn}
\setcounterpageref{ex}{0}

\begin{ex} %[2H5H26]
Tính độ dài đường gấp khúc gồm 2 đoạn thẳng có độ dài là $8\text{ cm}$ và $7\text{ cm}$ theo xăng-ti-mét.
\shortans{15}
\loigiai{
\textbf{Đáp số:} $15$\\
$8 + 7 = 15\text{ cm}$.
}
\end{ex}

\begin{ex} %[2H5H26]
Đường gấp khúc $ABC$ có $AB = 14\text{ cm}, BC = 16\text{ cm}$. Tính độ dài đường gấp khúc $ABC$ theo xăng-ti-mét.
\shortans{30}
\loigiai{
\textbf{Đáp số:} $30$\\
$14 + 16 = 30\text{ cm}$.
}
\end{ex}

\begin{ex} %[2H5V26]
Tính tổng độ dài các cạnh của một hình tứ giác có độ dài các cạnh là $5\text{ cm}, 6\text{ cm}, 7\text{ cm}, 8\text{ cm}$ theo xăng-ti-mét.
\shortans{26}
\loigiai{
\textbf{Đáp số:} $26$\\
$5 + 6 + 7 + 8 = 26\text{ cm}$.
}
\end{ex}

\begin{ex} %[2H5V26]
Đường gấp khúc $ABCD$ dài $35\text{ cm}$. Biết tổng độ dài đoạn $AB$ và $BC$ là $20\text{ cm}$. Tính độ dài đoạn $CD$ theo xăng-ti-mét.
\shortans{15}
\loigiai{
\textbf{Đáp số:} $15$\\
$CD = 35 - 20 = 15\text{ cm}$.
}
\end{ex}

\begin{ex} %[2H5V26]
Đường gấp khúc gồm 3 đoạn thẳng bằng nhau, mỗi đoạn dài $6\text{ cm}$. Tính độ dài đường gấp khúc đó theo xăng-ti-mét.
\shortans{18}
\loigiai{
\textbf{Đáp số:} $18$\\
$6 \times 3 = 18\text{ cm}$.
}
\end{ex}

\begin{ex} %[2H5C26]
Một hình tứ giác có tổng độ dài các cạnh là $40\text{ cm}$. Biết 3 cạnh có độ dài lần lượt là $8\text{ cm}, 10\text{ cm}, 12\text{ cm}$. Tính độ dài cạnh còn lại theo xăng-ti-mét.
\shortans{10}
\loigiai{
\textbf{Đáp số:} $10$\\
Cạnh còn lại: $40 - (8 + 10 + 12) = 40 - 30 = 10\text{ cm}$.
}
\end{ex}
